\section{Related Work}
Chimera detection in NGS has been widely studied in amplicon sequencing, where tools such as UCHIME and UCHIME2 test whether a query sequence can be explained as a mosaic of putative parents under reference-based or de novo assumptions. These approaches are effective in many microbial workflows but depend on database completeness or abundance hierarchies that do not reliably hold in single-species organellar sequencing and assembly.

Ensemble approaches such as CATCh combine multiple chimera detectors using supervised learning, improving performance in their target domain but inheriting the limitations and feature coverage of the underlying tools. Pipelines such as ChimPipe address biological chimeras in RNA-seq by integrating discordant paired-end and split-read signals, but are designed for fusion detection and typically assume available genome/annotation resources.

In contrast, mitochondrial assembly pipelines generally do not provide explicit read-level chimera filtering tailored to PCR artifacts. This gap motivates a reference-light, feature-driven approach that targets alignment fragmentation and compositional discontinuity in mitochondrial paired-end reads and evaluates success in terms of downstream assembly graph outcomes.